\documentclass[11pt,a4paper]{article}
\usepackage[T1]{fontenc}
\usepackage{isabelle,isabellesym}
\usepackage{amsmath,amssymb}
\usepackage{url}
\usepackage[pass]{geometry}
\usepackage{pdfsetup}
\isabellestyle{it}
\title{The Rank-One Dixmier Conjecture for Elements of Mass at Most Six}
\author{Idris Ali Shaik}
\date{}
\begin{document}
\maketitle
\begin{abstract}
This development formalizes a generation theorem for the first Weyl algebra
 over every field of characteristic zero. If $QP-PQ=1$ and $P$ has at most
six nonzero homogeneous components for the grading $\deg X=1$, $\deg Y=-1$,
then $P$ and $Q$ generate the algebra. No bound is imposed on the degree of
either member or on the mass of $Q$. The argument follows the mass-six
paper and its registered Lean formalization, including the structural
Newton and companion arguments used in the proof.
\end{abstract}

\section{Statement and representation}
For a field $K$ of characteristic zero, write
\[
 A_1(K)=K\langle X,Y\rangle/(YX-XY-1),\qquad [U,V]=UV-VU.
\]
The rank-one Dixmier conjecture asks whether every pair with $[Q,P]=1$
generates $A_1(K)$ \cite{Dixmier1968}. The result treated here has the
additional hypothesis $m(P)\leq 6$ \cite{Shaik2026Paper}.

The formalization represents $X$ by multiplication by the polynomial
variable and $Y$ by differentiation on $K[t]$. Thus $YX-XY$ is the identity
operator. The carrier \texttt{weyl\_algebra} is the unital $K$-subalgebra of
polynomial operators generated by these two operators. The PBW development
recovers their finite normal form $\sum_{i,j}c_{ij}X^iY^j$ and its commuting
symbol $\sum_{i,j}c_{ij}x^iy^j$.

For $P=\sum_g P_g$ in the $\mathbb Z$-grading above, mass is
\[
 m(P)=\#\{g:P_g\ne 0\}
     =\#\{i-j:c_{ij}\ne 0\}.
\]
In particular, mass counts occupied grades, not individual PBW monomials;
a single component can contain arbitrarily many monomials. The native
\texttt{weyl\_mass} definition counts the image of the finite symbol support
under $(i,j)\mapsto i-j$.

The public entry declaration is
\texttt{Dixmier\_\allowbreak Mass\_\allowbreak Six.\allowbreak mass\_\allowbreak six\_\allowbreak generation}. It is a direct alias of
the source producer
\texttt{Mass\_\allowbreak Six\_\allowbreak Generation\_\allowbreak Proved.\allowbreak massSixGeneration}.
Its field parameter has Isabelle sort \texttt{field\_char\_0}. It states
\[
 \begin{aligned}
 &P,Q\in\texttt{weyl\_algebra},\\
 &Q\circ P-P\circ Q=\mathrm{id},\qquad
   \texttt{weyl\_mass}\ P\leq 6\\
 &\quad\Longrightarrow\quad
   \texttt{op\_adjoin}\ \{P,Q\}=\texttt{weyl\_algebra}.
 \end{aligned}
\]
Here \texttt{op\_adjoin} is closure under the unital $K$-algebra operations,
so the conclusion is generation of the entire operator carrier. There is
no extra hypothesis that a structural input holds. The complex theorem is
\texttt{massSixGeneration\_complex}; scalar-carrier descent gives the stated
result over arbitrary characteristic-zero fields.

\section{Proof route}
The proof follows the existing architecture in
\cite{Shaik2026Paper,Shaik2026Lean}. Newton leading forms and homogeneous
companions reduce a nongenerating pair to a crossing case, a case already
forcing mass at least ten, or a total symbol with two distinguished roots.
Fourier transport preserves mass and permits the case reduction to be
applied to either orientation. The relevant structural arguments from
Guccione, Guccione and Valqui and Han and Tan are developed inside the
formalization \cite{GGV2014,GGV2025,HanTan2024}.

For a strict negative crossing, the companion identity becomes an equation
in complex univariate polynomials. The mass bound restricts the number of
nonzero coefficients of a power. Sparse-root and degree arguments force
the scalar polynomial into a quadratic binomial form, which is excluded by
the pure-power Newton-face argument. The horizontal crossing has its own
scalar reduction and leads to the same exclusion. In the two-root case,
root-multiplicity bounds on sparse polynomials contradict the degree
restrictions of a counterexample. These exclusions close the complex case;
the carrier and coefficient-field transfer closes the general theorem.

The development supplies the six structural inputs through their actual
proofs. The final theorem calls the proved degree-bound route and requires
only the carrier, commutator and mass hypotheses displayed above. It does
not assume a degree bound or a mass bound for the mate. The proof concerns
the mass-six theorem, and makes no assertion that the unrestricted
rank-one conjecture has been settled.

\section{Source and assistance}
The mathematical source is Theorem 1.1 of the mass-six paper, version 1.10.6
\cite{Shaik2026Paper}. The registered Lean result is
\texttt{Dixmier.\allowbreak Palomar.\allowbreak massSixGeneration}, Palomar entry
PALOMAR-2026-10-05-000006, version 2 \cite{Shaik2026Lean}. Its pinned source
commit is \nolinkurl{61783d52b6ae44cd2d8d20ad6cb798e7bbbce3ff}.
The Isabelle development reuses that proof architecture; its proof object
is the native Isabelle theory development below.

The mathematical conception, manuscript, verification of the mathematical
results, and part of the Lean 4 formalization were carried out by the author,
with AI assistance as disclosed in the respective manuscript or code.
OpenAI Codex generated a substantial portion of the Isabelle/HOL formalization
while porting the existing Lean 4 development, including proof reconstruction,
supporting lemmas, proof-script adaptation, and build troubleshooting.
The resulting formal proofs were checked by Isabelle's kernel.

\bibliographystyle{abbrv}
\bibliography{root}
\clearpage
\newgeometry{margin=18mm}
\renewcommand{\isastyle}{\normalfont\footnotesize\itshape}
\renewcommand{\isastyletext}{\small\normalfont\rmfamily}
\sloppy
\def\ {\space\allowbreak}
\section{Formal development}
\input{session}
\end{document}
